\section{Free-product geometry and invariant tensors}
\label{sec:tensors}

Let $D$ be a ${\rm II}_1$ factor generated as a von Neumann algebra
by an increasing sequence
$(D_j)$ of unital full matrix subalgebras whose sizes tend to infinity.
Let $A_0$ be a finite von Neumann algebra with a faithful normal tracial
state, and form the tracial von Neumann free product
\[
  \mathcal L=D*A_0.
\]
We use $\tau$ for the free-product trace and its restrictions.  Freeness
here is scalar tracial freeness: the trace of an alternating product of
centered letters from the two algebras is zero.  Its Hilbert-space
consequence is the orthogonal decomposition into reduced words, with
each reduced-word space isometric to the corresponding Hilbert tensor
product. This model comes from the free-product framework of
\citet{Ching1973} and \citet{Voiculescu1985}. The corner estimate below
is proved here from that word decomposition.

This section and the next two prove the following estimate.  Its constant
does not depend on the trace of the corner.

\begin{theorem}[Free-product corner estimate]\label{thm:free-corner}
Let $\mathcal L=D*A_0$ be as above, let $t=1/N$ for an integer $N\geq2$,
and suppose that $G\in\mathcal B(L^2(\mathcal L))$ commutes with $L_D$.
Suppose that $\epsilon\geq0$ and
\begin{equation}\label{eq:free-corner-hypothesis}
  \|[G,L_x]\|\leq\epsilon\|x\|
  \qquad
  (x\in e\mathcal L e,
   \ e\in D\text{ a projection with }\tau(e)=t).
\end{equation}
If $s=s^*=s^{-1}\in A_0$ and $\tau(s)=0$, then
\begin{equation}\label{eq:free-corner-conclusion}
 \left\|P_{L^2(A_0)}[G,L_s]\big|_{L^2(A_0)}\right\|
       \leq c_1\epsilon,
 \qquad c_1=6\pi^2(1+4c_0).
\end{equation}
\end{theorem}

Since equal-trace projections in $D$ are unitarily conjugate,
commutation with $L_D$ makes it equivalent to assume
\eqref{eq:free-corner-hypothesis} for one such projection.  Moreover,
for $x\in e\mathcal L e$ the commutator in that equation is supported on
$eL^2(\mathcal L)$.  The proof of the theorem is completed after the
Fourier compression calculation in Section~\ref{sec:hankel}.

\subsection{Factoriality and symmetrization}

\begin{lemma}\label{lem:free-factor}
One has $D'\cap\mathcal L=\mathbb C1$.  In particular, $\mathcal L$ is a
${\rm II}_1$ factor.
\end{lemma}

\begin{proof}
Put $E=L^2(D)$, $E^\circ=E\ominus\mathbb C1$, and
$B^\circ=L^2(A_0)\ominus\mathbb C1$.  Regrouping reduced words by their
first and last $A_0$-letters gives the $D$--$D$ bimodule decomposition
\[
 L^2(\mathcal L)\ominus L^2(D)
 \cong E\otimes\mathcal V\otimes E,
 \qquad
 \mathcal V=
 \bigoplus_{n\geq1}
 B^\circ\otimes(E^\circ\otimes B^\circ)^{\otimes(n-1)}.
\]
The endpoints in $E$ are unrestricted; expanding each as
$\mathbb C1\oplus E^\circ$ recovers every reduced word containing an
$A_0$-letter exactly once.  Left and right $D$ act on the first and last
factors, respectively.

Choose a Haar unitary $v$ in a diffuse abelian subalgebra of $D$ and put
$u_n=v^n$.  These unitaries converge ultraweakly to zero in $D$.
Indeed, the restriction of any normal functional to
$W^*(v)\cong L^\infty(\mathbb T)$ has an $L^1(\mathbb T)$ density,
whose Fourier coefficients tend to zero.  Consequently
$L_{u_n}$ and $R_{u_n^*}$ converge weakly to zero on $E$: first test
against bounded vectors in $D$, and then use their $L^2$ density.
On elementary tensors the matrix coefficients of
\[
 L_{u_n}R_{u_n^*}
  =L_{u_n}\otimes I_{\mathcal V}\otimes R_{u_n^*}
\]
therefore tend to zero.  The same holds on the whole complement by
density and uniform boundedness.  A $D$-central vector in this
complement is fixed by all these operators, and hence is zero.

If $z\in D'\cap\mathcal L$, its component perpendicular to $L^2(D)$
is $D$-central, because the displayed decomposition is a bimodule
decomposition.  Thus $z=E_D(z)\in D$, where $E_D$ is the
trace-preserving conditional expectation.  It follows that
$z\in Z(D)=\mathbb C1$.  The center of $\mathcal L$ is contained in
this relative commutant.  Its faithful finite trace and its diffuse
subfactor $D$ then make $\mathcal L$ a ${\rm II}_1$ factor.
\end{proof}

For $u\in\mathcal U(D)$, let $\theta_u$ be the trace-preserving
free-product automorphism which equals $\operatorname{Ad}u$ on $D$
and is the identity on $A_0$.  Its Hilbert-space implementation $W_u$
acts by $\theta_u$ on bounded words.  In particular,
\[
 W_uL_xW_u^*=L_{\theta_u(x)},\qquad
 W_u\xi=\xi\quad(\xi\in L^2(A_0)).
\]

\begin{lemma}[Symmetrization]\label{lem:free-symmetrization}
In proving Theorem~\ref{thm:free-corner}, one may suppose in addition
that $G$ commutes with $W_u$ for every
$u\in\bigcup_j\mathcal U(D_j)$.  This replacement preserves the
compressed commutator in \eqref{eq:free-corner-conclusion} exactly,
and preserves \eqref{eq:free-corner-hypothesis} for both $G$ and $G^*$.
\end{lemma}

\begin{proof}
Average the original operator over each compact matrix unitary group:
\[
 G_j=\int_{\mathcal U(D_j)}W_uGW_u^*\,d\mu_j(u).
\]
This weak operator integral is well defined.  The unitary
representation $u\mapsto W_u$ is strongly continuous on each group,
as is seen first on bounded reduced words and then by density.
Every conjugate in the integral commutes with $L_D$.  It also obeys
\eqref{eq:free-corner-hypothesis}, since $\theta_u$ permutes the
$D$-corners of trace $t$.  Convexity and weak operator closedness of
norm balls preserve these properties under integration and under a
weak operator cluster limit $\widetilde G$ as $j\to\infty$.
Such a limit exists because $\|G_j\|\leq\|G\|$.
For fixed $u\in\mathcal U(D_i)$, all $G_j$ with $j\geq i$ commute
with $W_u$, by Haar invariance.  Thus $\widetilde G$ has every stated
symmetry.

The operators $W_u$ fix $L^2(A_0)$ pointwise and commute with $L_s$.
For each conjugate, and therefore for the averages and their limit,
\[
 P_{L^2(A_0)}[W_uGW_u^*,L_s]\big|_{L^2(A_0)}
   =P_{L^2(A_0)}[G,L_s]\big|_{L^2(A_0)}.
\]
Finally,
$[G^*,L_x]=-[G,L_{x^*}]^*$, so the hypothesis for a self-adjoint
corner algebra also gives its adjoint version.  The same applies to
$\widetilde G$.
\end{proof}

For the remainder of the proof of Theorem~\ref{thm:free-corner},
$G$ denotes the symmetrized operator.  At the end we recover the
original compressed commutator using
Lemma~\ref{lem:free-symmetrization}.

\subsection{The word decomposition}

Set
\[
 E=L^2(D),\quad E^\circ=E\ominus\mathbb C1,
 \qquad B=L^2(A_0),\quad B^\circ=B\ominus\mathbb C1,
 \qquad J_k=(B^\circ)^{\otimes k}\otimes B.
\]
We use the following ordering of the reduced-word factors:
\begin{equation}\label{eq:word-decomposition}
 L^2(\mathcal L)
   =\bigoplus_{k\geq0}\mathcal H_k,
 \qquad
 \mathcal H_k=E\otimes(E^\circ)^{\otimes k}\otimes J_k.
\end{equation}
The tensor
\[
 d_0\otimes d_1\otimes\cdots\otimes d_k
       \otimes(a_1\otimes\cdots\otimes a_k\otimes z)
\]
represents the word $d_0a_1d_1\cdots a_kd_kz$.  Initially the
letters may be taken bounded, with $d_i\in D\ominus\mathbb C1$
and $a_i\in A_0\ominus\mathbb C1$ for $i\geq1$, and unrestricted
$d_0\in D$, $z\in A_0$.  Expanding the endpoints into scalar and
centered parts produces all the usual reduced-word spaces exactly
once.  This proves both orthogonality and completeness in
\eqref{eq:word-decomposition}; the tensor notation then extends to
arbitrary $L^2$-letters by the corresponding isometries.
On each level, $L_D$ acts on the first factor, while $W_u$ conjugates
all $k+1$ factors coming from $D$ and fixes $J_k$.

We will show that $G$ preserves each word level $\mathcal H_k$ and, on
tensors with a repeated centered $D$ letter, changes only the $J_k$ factor.
Scalar slices against fixed vectors in the $J_k$ and $J_l$ factors reduce
these assertions to the following intertwiner calculation.
Write $C_u\xi=u\xi u^*$ on $E$ and on its invariant
subspace $E^\circ$.

\begin{lemma}[Invariant intertwiners]\label{lem:invariant-intertwiners}
Let $k,l\geq0$, and let
\[
 Y:E\otimes(E^\circ)^{\otimes k}
       \longrightarrow E\otimes(E^\circ)^{\otimes l}
\]
be bounded.  Suppose that $Y$ intertwines the left $D$ actions and
that
\[
 Y C_u^{\otimes(k+1)}=C_u^{\otimes(l+1)}Y
 \qquad(u\in\bigcup_j\mathcal U(D_j)).
\]
Then $Y=0$ if $k\ne l$.  If $k=l$, there are scalars
$(a_\sigma)_{\sigma\in S_k}$ such that
\[
 Y=I_E\otimes\sum_{\sigma\in S_k}a_\sigma U_\sigma,
 \qquad
 U_\sigma(x_1\otimes\cdots\otimes x_k)
   =x_{\sigma^{-1}(1)}\otimes\cdots\otimes x_{\sigma^{-1}(k)}.
\]
For $k=0$ the sum consists of the scalar multiple of the identity
corresponding to the empty permutation.
\end{lemma}

\begin{proof}
For a vector or algebra element $x\in E$, put
$x^\circ=x-\tau(x)1$.  On unrestricted variables
$a,x_1,\ldots,x_k,z_1,\ldots,z_l\in D$, define the complex
multilinear form
\[
 F(a,x,z)=
 \left\langle
 Y(1\otimes x_1^\circ\otimes\cdots\otimes x_k^\circ),
 a^*\otimes(z_1^*)^\circ\otimes\cdots\otimes(z_l^*)^\circ
 \right\rangle.
\]
It is multilinear with our convention that inner products are linear
in the first argument.  Fixing the tail vectors and taking their
Hilbert-space slice of $Y$ gives an operator $T:E\to E$ commuting
with $L_D$.  The tracial commutation theorem gives $T=R_b$ for a
bounded $b\in D$, with
\[
 \|b\|=\|T\|
  \leq\|Y\|\prod_i\|x_i^\circ\|_2
                 \prod_j\|(z_j^*)^\circ\|_2.
\]
Since $F(a,x,z)=\tau(ab)$, this yields the endpoint estimate
\begin{equation}\label{eq:coefficient-bound}
 |F(a,x,z)|\leq
 \|Y\|\,\|a\|_1
       \prod_{i=1}^k\|x_i\|_2
       \prod_{j=1}^l\|z_j\|_2.
\end{equation}
The same slice also identifies coefficients with arbitrary bounded
first endpoints: replacing $1,a^*$ by $d,c$ replaces $a$ by $c^*d$,
because $\langle db,c\rangle=\tau(c^*db)$.

\smallskip\noindent
\emph{Finite matrix invariants and stable coefficients.}
Put $p=1+k+l$ and restrict all variables to a stage
$D_j\cong M_m(\mathbb C)$.  The intertwining assumption says that
$F$ is invariant under simultaneous unitary conjugation.  To see
precisely the finite-dimensional consequence, label its variables
$v_0,\ldots,v_{p-1}$.  Every multilinear form on $M_m^p$ has a unique
expression
\[
 \operatorname{Tr}\bigl(A(v_0\otimes\cdots\otimes v_{p-1})\bigr)
 \qquad
 (A\in\operatorname{End}((\mathbb C^m)^{\otimes p})),
\]
by the nondegenerate trace pairing.  Conjugation invariance is
equivalent to $A$ commuting with every $u^{\otimes p}$.
Finite-dimensional Schur--Weyl duality identifies this commutant
with the span of the tensor permutation operators; they are linearly
independent when $m\geq p$
\citep[Theorem~3.1]{GrinkoOzols2024}.
Taking their trace pairings gives the cycle-trace forms
\[
 \Phi_\pi(v_0,\ldots,v_{p-1})
   =\prod_{(i_1\,\cdots\,i_r)\text{ a cycle of }\pi}
       \tau(v_{i_1}\cdots v_{i_r}),
 \qquad \pi\in S_p.
\]
Here cycles follow the successor convention of the displayed
product; reversing the convention only reindexes the permutations.
The normalization from $\operatorname{Tr}$ to $\tau$ multiplies
each form by a nonzero scalar, so these forms are also independent
for $m\geq p$.  Thus
\[
 F\big|_{D_j^p}=\sum_{\pi\in S_p}\lambda_{\pi,j}\Phi_\pi.
\]
This identity is on \emph{unrestricted} matrix variables: the
centering occurs inside the definition of $F$.
Under a unital inclusion of full matrix stages, normalized traces
agree, and hence every $\Phi_\pi$ restricts to the identically
defined form on the smaller stage.  Uniqueness on any smaller stage
of size at least $p$ proves that $\lambda_{\pi,j}$ is independent
of $j$ from that point onward.  Denote it by $\lambda_\pi$.

\smallskip\noindent
\emph{Isolating coefficients before restricting to centered tails.}
The invariant-theory step has expressed every coefficient through
finitely many cycle-trace forms with coefficients fixed across matrix
stages. We now use the endpoint bound to remove long cycles, then
Hilbert-space boundedness to remove pairs within one side.
Fix $\pi\in S_p$.  Choose distinct matrix indices for the slots and
put $v_i=e_{i,\pi(i)}$.  In a nonzero product of these matrix units,
the variable in slot $i$ must be followed by the variable whose
initial index is $\pi(i)$.  Closing each product under its trace
therefore forces exactly the cycles of $\pi$.  Consequently
\[
 \Phi_\rho(v_0,\ldots,v_{p-1})
  =\begin{cases}
      m^{-c(\pi)},&\rho=\pi,\\
      0,&\rho\ne\pi,
    \end{cases}
\]
where $c(\pi)$ is the number of cycles.  Every chosen matrix unit
has $L^1$ norm $m^{-1}$ and $L^2$ norm $m^{-1/2}$, including
diagonal units.  Applying \eqref{eq:coefficient-bound} gives
\[
 |\lambda_\pi|m^{-c(\pi)}
       \leq\|Y\|m^{-1-(k+l)/2}.
\]
The coefficient is fixed while the stage size tends to infinity.
It follows that
\[
 \lambda_\pi=0\qquad\text{if}\qquad
 c(\pi)<1+\frac{k+l}{2}.
\]

Now restrict the $k+l$ tail variables to centered matrices.  Any
form with a singleton tail cycle vanishes.  For a permutation with
no such singleton, let $r$ be the length of the cycle containing the
distinguished variable $a=v_0$.  Its number of cycles is at most
\[
 1+\frac{k+l-(r-1)}2.
\]
If $r>1$, this is below the required threshold.  If $r=1$, the
threshold is attained only when every tail cycle has length two.
Thus the only remaining terms have $a$ as a singleton and pair all
tail slots.  On centered tails,
\[
 F(a,x,z)=\tau(a)
       \sum_{\mathcal P}\lambda_{\mathcal P}
          \prod_{\{i,j\}\in\mathcal P}\tau(v_i v_j),
\]
where $\mathcal P$ ranges over pairings of the tail slots.  There
are no terms when the number of tail slots is odd.

Using $a=c^*d$ gives the corresponding simple-tensor coefficients
with first endpoints $d,c$ in the union of the matrix stages.
These formulas extend to arbitrary endpoints in $E$ and arbitrary
centered $L^2$ tail letters.  Indeed, that union is $L^2$ dense,
the left side is continuous in its individual $L^2$ letters, and
each factor on the right is continuous by
$|\tau(uv)|\leq\|u\|_2\|v\|_2$.  At this point the assertion is
continuity of coefficients on simple tensors, not boundedness of
each pairing on the completed tensor product.

\smallskip\noindent
\emph{Excluding pairings within one side.}
Choose an infinite orthonormal family $(h_n)$ of self-adjoint
centered vectors in $E$.  For example, normalized real
trigonometric functions of a Haar unitary give such a family.
Then $\tau(h_nh_m)=\delta_{nm}$.
Fix a pairing $\mathcal P$, with $r_s$ pairs entirely in the source
slots and $r_t$ entirely in the test slots.  Assign a disjoint
infinite pool of the $h_n$ to each pair.  For each source--test pair,
use one fixed member of its pool in both slots.  For each
source--source pair, sum the source tensor over $M$ members of its
pool, placing the same member in its two slots, independently for
the different pairs.  Form the test tensor analogously for the
test--test pairs.  Take both first endpoints to be $1$.

In any term of the expanded coefficient every label occurs exactly
twice, in the positions assigned by $\mathcal P$, and different
pairs have different labels.  Hence all other pairings contribute
zero.  The resulting coefficient is exactly
$\lambda_{\mathcal P}M^{r_s+r_t}$.  The source summands are
orthonormal, as are the test summands, so the product of the two
tensor norms is $M^{(r_s+r_t)/2}$.  Boundedness of $Y$ gives
\[
 |\lambda_{\mathcal P}|M^{r_s+r_t}
       \leq\|Y\|M^{(r_s+r_t)/2}.
\]
Letting $M\to\infty$ proves $\lambda_{\mathcal P}=0$ whenever
$r_s+r_t>0$.  The argument also applies when one side has no tail
slots.

All surviving pairs therefore join a source slot to a test slot,
which requires $k=l$.  If these numbers differ, all simple-tensor
coefficients vanish and $Y=0$.  Otherwise these pairings are
indexed by bijections $\sigma\in S_k$, and the coefficients are
\[
 \tau(c^*d)\sum_{\sigma\in S_k}a_\sigma
       \prod_{i=1}^k\tau(x_i z_{\sigma(i)}).
\]
Since $\tau(x_i z_j)=\langle x_i,z_j^*\rangle$, these are precisely
the coefficients of
$I_E\otimes\sum_\sigma a_\sigma U_\sigma$.  Equality against all
elementary test tensors, and their density, proves the asserted
operator identity.
\end{proof}

\begin{lemma}[Word structure]\label{lem:word-structure}
The symmetrized operator $G$ preserves every level $\mathcal H_k$
in \eqref{eq:word-decomposition}.  There are operators
$B_k\in\mathcal B(J_k)$, with $\|B_k\|\leq\|G\|$, such that, for
all $d_0\in E$, $x\in E^\circ$, and $\alpha\in J_k$,
\begin{equation}\label{eq:repeated-letter}
 \begin{aligned}
 G(d_0\otimes x^{\otimes k}\otimes\alpha)
     &=d_0\otimes x^{\otimes k}\otimes B_k\alpha,\\
 G^*(d_0\otimes x^{\otimes k}\otimes\alpha)
     &=d_0\otimes x^{\otimes k}\otimes B_k^*\alpha.
 \end{aligned}
\end{equation}
These are identities in the full space $L^2(\mathcal L)$.
The operators $B_k$ are independent of the repeated letter $x$;
in particular, $x$ may be any complex centered $L^2$ vector.
\end{lemma}

\begin{proof}
Let $P_k$ be the level projection.  It commutes with $L_D$ and all
$W_u$.  For $\alpha\in J_k$ and $\beta\in J_l$, take the scalar
slice in these factors of $P_lGP_k$ to obtain an operator
\[
 Y_{\alpha,\beta}:
 E\otimes(E^\circ)^{\otimes k}
   \longrightarrow E\otimes(E^\circ)^{\otimes l},
 \qquad
 \|Y_{\alpha,\beta}\|\leq\|G\|\|\alpha\|\|\beta\|.
\]
It satisfies the hypotheses of
Lemma~\ref{lem:invariant-intertwiners}.  For $k\ne l$ every such
slice vanishes, so $P_lGP_k=0$.  Thus the levels reduce $G$.

Fix a unit vector $h\in E^\circ$ and let
$V_k:J_k\to\mathcal H_k$ be the isometry
$V_k\alpha=1\otimes h^{\otimes k}\otimes\alpha$.
Define $B_k=V_k^*GV_k$.  On an equal-level scalar slice, the lemma
gives
\[
 Y_{\alpha,\beta}
    =I_E\otimes\sum_{\sigma\in S_k}
                        a_\sigma(\alpha,\beta)U_\sigma,
 \qquad
 \langle B_k\alpha,\beta\rangle
       =\sum_{\sigma\in S_k}a_\sigma(\alpha,\beta).
\]
Every $U_\sigma$ fixes $x^{\otimes k}$ for every $x\in E^\circ$.
Consequently the last scalar is independent of the unit letter
used to define $V_k$.  The same slice identity, now tested against
arbitrary endpoint and tail vectors and arbitrary $\beta$, gives
the first equality in \eqref{eq:repeated-letter}.  These tests
determine the whole output on level $k$, and the other levels have
already been excluded.  This also proves the assertion for every
endpoint $d_0\in E$, with no boundedness requirement on its
representative.  No self-adjointness of $x$ is used: the relevant
cross contraction is $\tau(xz)=\langle x,z^*\rangle$.

The same argument applies to $G^*$.  Its induced operator on $J_k$
is $V_k^*G^*V_k=(V_k^*GV_k)^*=B_k^*$, proving the second equality.
\end{proof}

At level zero the lemma says $G|_{\mathcal H_0}=I_E\otimes B_0$.
Since $B=L^2(A_0)$ is the subspace $\mathbb C1\otimes B$ there,
the quantity to be estimated in
Theorem~\ref{thm:free-corner}, after symmetrization, is exactly
\[
 \left\|P_B[G,L_s]\big|_B\right\|
      =\|[B_0,L_s|_B]\|.
\]
